\section{The experimental problem}\label{sec:experiment}
The mother problem is to pass from a physically executable experiment to a quantitative law for the predictions that can actually be retained. There are two distinct obstructions. A set of formally reachable predictions need not receive appreciable mass under the preparation and exploration being used. Moreover, a good codebook at a checkpoint need not admit a good recursive implementation. The first obstruction is probabilistic; the second is causal. Treating either as a dimension count loses the mechanism of the problem.

Our main theorem connects these obstructions by weighting the propagation of quantization errors with the same acquisition probabilities that enter the lower bound. Inactive steps are especially important. When an update preserves a retained representative exactly, rounding need not be repeated. Charging an error at every physical tick would produce a spurious inverse acquisition probability in a system that simply waits. The theorem retains the actual clock while localizing error injection to the updates that require it.

The construction of states from predictions is established background; controlled predictive representations already use future action--observation tests \cite{psr}. Squared-loss projection and probability quantization are classical \cite{graf}. Recursive quantization of nonlinear filters has a substantial literature \cite{pages,sellami,feuer}. Experiment comparison originates in particular with Blackwell \cite{blackwell}. Our contribution here is not a claim of priority for any of these ingredients. It is the acquired-mass/innovation-transport theorem below, its all-budget rank-uniform geometric consequence, and the two kernel-level verifications. Neither realization is used in the proof of the general theorem.

\subsection{Raw data, strategies, and preparation}
A discrete-time experiment consists of standard Borel spaces $X_t,A_t,Y_{t+1}$, a probability preparation $\mu$ on $X_0$, and nonnegative normalized kernels
\begin{equation}\label{eq:raw}
 K_t(dx',dy,dc\mid x,a).
\end{equation}
The visible history $h_t$ records actions, reports, and the published nonnegative resource increments $c$. These increments include failed preparations and physical elapsed time. Positivity means nonnegativity and normalization, not a positive lower bound for every report. Deterministic transitions are permitted. An unknown physical parameter, when present, is included in the hidden state with a specified prior for the Bayesian assertions below. Parameter-uniform experiment comparison will be stated separately.

A legal strategy is a stochastic kernel $\alpha_t(da\mid h_t)$ supported on the declared available actions. All randomness used by the apparatus or algorithm is independent of the hidden preparation unless generated by \eqref{eq:raw}. A stopping rule is a stopping time for the visible filtration. The main risk theorem fixes a deterministic raw-acquisition horizon $N$ and an exploration strategy $\beta$ independently of the encoder. It does not optimize $\beta$. At most $N$ actual calls, including unsuccessful calls, are made. A stopped comparison is always explicitly bounded by $N$.

Successive kernel integration gives the path law $P^\alpha$. The complete conditional belief $\pi_t(h_t)$ on the hidden state has a measurable version. We specify versions on a declared admissible history domain; identities on impossible reports are not inferred by division by zero.

\subsection{Executable future tests and the quotient}
A future test is a legal finite continuation followed by a bounded measurable function of its actual reports. Hidden-state functions qualify only when a terminal measurement producing them is part of the experiment. Tests include determining report events and the continuation interface: legal actions, remaining budget, and every clock value affecting future experiments. They are closed under legal one-step continuation, finite randomization, and the bounded monotone limits needed to determine the designated laws.

For an admissible prepared belief $b$, let $\mathsf e_t(b)$ be its evaluations on a countable determining family of these tests, together with the continuation interface. Equality of $\mathsf e_t$ is \emph{predictive equivalence}. We use the following verifiable realization certificate. The admissible belief domains $\mathcal B_t$ and the action domains are compact metric; the determining evaluations and the report probabilities are continuous. For continuous reports on a declared compact report domain, there is a common dominating measure and a jointly continuous normalized update on the declared domain. Where the Bayes ratio is used its density denominator is strictly positive. Any extensions at zero denominators are separately verified to be continuous and predictive-fibre-compatible. Alternatively, finite reports with continuous normalized updates suffice. The update is assumed to preserve the declared belief domains. The test family determines the full designated continuation laws. Compactness is a certificate for this class, not an axiom asserting that every experiment has a compact quotient.

\begin{proposition}[Predictive realization and minimality]\label{prop:quotient}
Under this certificate, $S_t=\mathsf e_t(\mathcal B_t)$ is compact metrizable. It realizes predictive equivalence, and the report probabilities and the normalized update descend to continuous functions and maps on $S_t$, on the domain where the latter is specified. Any measurable history statistic $r_t$ with respect to which every determining test evaluation and the continuation interface are measurable admits a measurable factor $s_t=f_t(r_t)$ almost surely. A set-theoretic sufficient statistic admits the analogous unique factor on its image, without an automatic Borel assertion.
\end{proposition}
\begin{proof}
Rescale evaluations into $[-1,1]$. The countable product is compact metrizable and $\mathsf e_t$ is continuous. Its compact image is therefore compact metrizable. Equality on the determining tests implies equality of all designated future laws by induction on continuation depth and a monotone-class argument. In the induction step the numerator and denominator of each conditional test are evaluations of one-step pullbacks; hence equivalent beliefs give equivalent updated beliefs whenever the denominator is positive.

A continuous surjection from a compact space to a Hausdorff space is a quotient map. The same assertion applies after taking the product with the compact action domain and, for continuous reports, the declared compact report domain. Thus every continuous fibre-constant evaluation or update descends continuously. The finite-report case is applied separately to each report. For measurable minimality, apply the measurable factorization lemma to each real evaluation as a $\sigma(r_t)$-measurable random variable, then combine the countably many factors; on the full-probability set where the resulting vector belongs to $S_t$, it is the required factor. Outside that set assign a fixed point of $S_t$. Setwise minimality follows directly from inclusion of the fibres of $r_t$ in predictive fibres.
\end{proof}

Finite exact continuation-interface components can be included as a disjoint compact union. Noncompact report spaces require an independent Borel realization or a verified localization; no universal smooth or standard Borel quotient of an arbitrary equivalence relation is asserted. In the realizations below the quotient is computed explicitly from the raw kernel, which also fixes the conditional versions.

\subsection{The common prediction task}
At time $N$, after the label has been retained, an independent index $i\in\{1,\ldots,q\}$ is selected with fixed positive weights $v_i$, $\sum_i v_i=1$. The corresponding legal terminal probe returns $T_i\in[0,1]$. The algorithm reports $a_i\in[0,1]$ and incurs $(T_i-a_i)^2$. Each expectation refers to its own executable probe; a joint counterfactual vector of all probe outcomes is unnecessary. Put
\begin{equation}\label{eq:score}
 p_t(s)=\big(\sqrt{v_i}\,\E[T_i\mid s]\big)_{i=1}^q,
 \qquad B_N=\sum_i v_i\E\Var(T_i\mid H_N),
 \qquad \nu_N=\law(p_N(s_N)).
\end{equation}
A varying terminal time uses the corresponding declared probe. The finite scoring menu need not determine the entire predictive quotient. Our assumptions below require only its Lipschitz upper bound globally and a separating lower bound on each acquired chart. Unmeasured interface coordinates cannot be silently charged to this task's lower bound.

A checkpoint encoder reads $H_N$ and retains at most $m$ labels. An online encoder has at most $m$ persistent labels at every time, updates from its previous label and the current input/report, and has no readable history tape. A terminal decoder may depend on the public horizon and on the freshly selected probe index, but not on discarded observations. Shared randomness does not carry observation-dependent state. The risks $R_\cp(N,m)$ and $R_\on(N,m)$ are infima of the full expected score over these respective classes, under the fixed exploration $\beta$. Temporary workspace, numerical input access, and program description are recorded separately in Section~\ref{sec:transfer}.
